\cotefile{29}
\scanpage{22}
% archive sheet number: 21
faisceau associé au préfaisceau précédent (NB la condition canularesque mise sur les $(D_i)$ assure que l'on a bien un préfaisceau). Alors localement pour la topologie étale, on a donc une situation $(D_i)$ comme dans a). Cela suffit pour définir une donnée de ramification $\underline{R}$ sur $B$, généralisant d'ailleurs celle envisagée dans a) (du moins lorsque la famille $(D_i)_{i\quad I}$ de a) satisfait à la condition supplémentaire locale qu'on\handwritten{\unsure{g|?}} vient de formuler). Je vois deux cas particuliers ayant de l'intérêt;. \deleted{\unsure{S|R|fi|si}} et \deleted{1\textsuperscript{\uline{o}})} où $\quad$ peut être défini en termes d'un diviseur $D$ donné sur $S$:

1\textsuperscript{\uline{o}}) $S$ est localement noethérien, $D$ est à croisements normaux, on prend comme $D_i$ les composantes irréductibles locales régulières de $D$ (qui ne sont définies que localement pour la topologie étale).

2\textsuperscript{\uline{o}}) $S$ est au dessus d'un préschéma $T$, $D$ est à croisements normaux relativement à $T$, on prend comme $D_i$ les composantes locales lisses sur $\unsure{\deleted{T}S|\deleted{S}T}$ (qui ne sont définies que localement pour la topologie étale).

Dans le cas 1\textsuperscript{\uline{o}}), la catégorie $\operatorname{Rev}(_{\underline{R}}(S)$ est équivalente à une sous-catégorie pleine de la catégorie des revêtements de $S$, et sauf erreur (il faudrait le vérifier) il en est de même dans le cas 2\textsuperscript{\uline{o}}). Les revêtements en question s'appellent, dans le\deleted{ur} cas 1\textsuperscript{\uline{o}}), revêtements modérément ramifiés au sens large relativement au diviseur $D$ (NB on vérifie en effet que c'est une notion plus large que celle de ``revêtement modérément ramifié par rapport à $D$'', (signifiant : par rapport à la famille réduite à $D$). Le théorème d'Abhyankar permet de caractériser ces revêtements par la condition d'être étales sur \deleted{$S\quad\unsure{\underline{r}|\underline{a}|\underline{R}}$} $S-\operatorname{supp}D$, normaux en les points au dessus de $\operatorname{supp}D$, toute composante irréductible en d\unsure{\deleted{i}o|\deleted{o}i}mine une de $S$, enfin ils sont modérément ramifiés au sens classique en le\inserted{s} points maximaux de $D$. Dans le cas 2\textsuperscript{\uline{o}}, il y aurait lieu d'appeler les revêtements obtenus ``modérément ramifiés pour $D$, relativement à la
